% Appendix C: structural results. Ported 2026-11-21 from
% docs/paper4_continuation/CAUSAL_SAFETY_PREMIUM_THEORY.md (Lemma R, Z1, N1),
% theory_v2/STOCHASTIC_POSITIVE_PREMIUM.md (Lemmas S, Pre; SP1 at k = 0),
% theory_v2/MIXING_PREMIUM_THEOREM.md (M1, M3, M4, M10) and theory_v2/B_GENERAL_FIXED_PREMIUM_PROOF.md (M11).
\section{Structural results: proofs of Proposition~\ref{prop:z1} and Theorem~\ref{thm:premium}}
\label{app:structure}

The results in this appendix hold for any label delay $d\ge1$ unless stated. Throughout, a \emph{pure plan} is a deterministic pre-promotion strategy; after the first promotion the learner plays the optimal belief DP, which may roll back, re-promote and use any remaining labels. We write $\mathrm{gain}(\pi)=C_0-C(\pi)$ for the expected cost saving of a plan under the alternative and $R_{\mathrm{IID}}(\pi)$ for its promotion probability under the i.i.d.\ null.

\subsection{Lagrangian reduction and monotone promotion events}

Let $\mathcal P$ be the finite set of pure plans and $\mathcal M\subseteq\mathcal P$ the plans whose promotion event, as a function $g_\pi(y_1,\dots,y_T)$ of the full potential label tape, is coordinatewise nondecreasing. Never-promote belongs to $\mathcal M$.

\begin{lemma}[Lagrangian reduction]
\label{lem:R}
If $\min_{\pi\in\mathcal P}[C(\pi)+\lambda R_{\mathrm{IID}}(\pi)]=\min_{\pi\in\mathcal M}[C(\pi)+\lambda R_{\mathrm{IID}}(\pi)]$ for every $\lambda\ge0$, then $\Gamma=0$.
\end{lemma}

\begin{proof}
By Kuhn's theorem every strategy is a mixture of pure plans, and $C$ and $R_{\mathrm{IID}}$ are linear in the mixture. So $\Vi=\min\{\E_\mu C:\E_\mu R_{\mathrm{IID}}\le\alpha,\ \mu\in\Delta(\mathcal P)\}$, and likewise $V_{\mathrm{IID}}(\mathcal M)$ over $\Delta(\mathcal M)$. Both are finite linear programmes with one inequality, and never-promote is feasible with risk 0, so strong duality gives $V_{\mathrm{IID}}(\mathcal S)=\max_{\lambda\ge0}[\min_{\pi\in\mathcal S}(C+\lambda R_{\mathrm{IID}})-\lambda\alpha]$ for $\mathcal S\in\{\mathcal P,\mathcal M\}$. The hypothesis makes the two dual functions equal, so $V_{\mathrm{IID}}(\mathcal M)=\Vi$.

Take any mixture $\mu$ on $\mathcal M$ and any adaptive null $\nu$. Couple the null tape $Y_t=\one\{U_t\le p_t\}$ with $Y^\ast_t=\one\{U_t\le1/2\}$ for i.i.d.\ uniform $U_t$. Since $p_t\le1/2$, $Y_t\le Y^\ast_t$ pathwise for any adapted $p_t$, including $p_t$ that react to public actions revealing which plan was drawn. Monotonicity gives $g_\pi(Y)\le g_\pi(Y^\ast)$, so $\PP_\nu(\text{promote})\le R_{\mathrm{IID}}(\mu)$. The i.i.d.-optimal mixture over $\mathcal M$ is therefore causal-feasible, and $\Vi\le\Vc\le V_{\mathrm{IID}}(\mathcal M)=\Vi$.
\end{proof}

\subsection{Proof of Proposition~\ref{prop:z1}}

Write $\rho=1-\eta_{01}-\eta_{10}\ge0$ for the eigenvalue of the transition matrix. Fix a pure plan $\pi$ with one label. If $\pi$ never queries, or promotes at some $s<\tau+d$ before its label is released, its promotion event is constant (empty or everything); $\tau$ is deterministic because nothing has been observed before it. Otherwise, after the release the plan is open-loop given the label value $y=y_\tau$ and promotes at $s_y\in\{\tau+d,\dots,T\}\cup\{\infty\}$. Its event $g(y)=\one\{s_{y_\tau}<\infty\}$ is non-monotone only if $s_0<\infty=s_1$.

For $y\in\{0,1\}$ and $s\ge\tau+d$ let $G_y(s)$ be the expected cost saved by ``promote at $s$, then continue optimally'' over ``never promote'', given $Y_\tau=y$. After the single query the budget is 0, so the continuation is open loop and
\[
G_y(s)=\max_{a:\,a_s=1}\ \sum_{t\ge s}a_t\,(2q_t(y)-1)-k\cdot\mathrm{sw}(a),\qquad q_t(y)=e_0+(e_1-e_0)\PP(S_t=1\mid Y_\tau=y).
\]
Because $e_1>e_0$, the likelihood ratio is monotone and $\PP(S_\tau=1\mid Y_\tau=1)\ge\PP(S_\tau=1\mid Y_\tau=0)$. Propagation gives $\PP(S_t=1\mid y)=\pi_\infty+(\PP(S_\tau=1\mid y)-\pi_\infty)\rho^{t-\tau}$ with $\rho\ge0$, so $q_t(1)\ge q_t(0)$ for all $t\ge\tau$. Each schedule's value is nondecreasing in $(q_t)$ because $a_t\ge0$, and a maximum of nondecreasing functions is nondecreasing, so $G_1(s)\ge G_0(s)$.

Let $\pi$ have $s_0<\infty=s_1$, write $K$ for the common cost of never promoting after the query and $P_y=\PP_A(Y_\tau=y)$. Then $C(\pi)=K-P_0G_0(s_0)$ and $R_{\mathrm{IID}}(\pi)=1/2$.
(a) If $G_0(s_0)\le0$, replace $\pi$ by the same plan with $s_0=\infty$: the event is empty, the cost does not increase and the risk falls to 0.
(b) If $G_0(s_0)>0$, replace $\pi$ by $\pi'$ with $s'_1=s_0$ and $s'_0=\infty$. Its event $\{y_\tau=1\}$ is monotone, with the same i.i.d.\ risk $1/2$, and cost $K-P_1G_1(s_0)$. Since $G_1\ge G_0>0$ and $P_1\ge1/2\ge P_0$ by assumption, $P_1G_1\ge P_0G_0$ and $C(\pi')\le C(\pi)$.
In both cases the Lagrangian $C+\lambda R_{\mathrm{IID}}$ does not increase for any $\lambda\ge0$. Applying this to a Lagrangian-optimal pure plan yields a monotone one, and Lemma~\ref{lem:R} gives $\Gamma=0$. \hfill$\square$

Persistence is used only for $G_1\ge G_0$, and the marginal condition only to compare $P_1$ with $P_0$. Delay, costs, rollback and re-promotion enter only through $G_y$.

\subsection{One label: an exact formula, and Theorem~\ref{thm:premium}(1)}

\begin{lemma}[Exact premium at one label]
\label{lem:N1}
Let $B=1$ and $\alpha\le1/2$, with any two-state alternative. For each query time $\tau$ with $\tau+d\le T$, let $g_{z0}(\tau)$, $g_{z1}(\tau)$ and $g_{\mathrm{both}}(\tau)$ be the largest cost reductions against $C_0$ (query cost included) of designs that query at $\tau$ and promote, at the best post-release time, only after $Y_\tau=0$, only after $Y_\tau=1$, or after both values. Let $g_{\mathrm{pre}}$ be the best reduction of a plan that promotes unconditionally (before any release, or with no query). With
\[
r_c=\max\big(0,\,g_{\mathrm{pre}},\,\max_\tau\{2g_{z1}(\tau),\,g_{z0}(\tau),\,g_{\mathrm{both}}(\tau)\}\big),\qquad r_{\mathrm{IID}}=\max\big(r_c,\,2\max_\tau g_{z0}(\tau)\big),
\]
we have $\Vi=C_0-\alpha r_{\mathrm{IID}}$, $\Vc=C_0-\alpha r_c$ and $\Gamma=\alpha\,(2\max_\tau g_{z0}(\tau)-r_c)^+$.
\end{lemma}

\begin{proof}
A pure plan with one label has a constant event (risk 0 or 1 under every null), a $\{Y_\tau=1\}$ event (i.i.d.\ and adaptive risk $1/2$), or a $\{Y_\tau=0\}$ event (i.i.d.\ risk $1/2$, adaptive risk 1). At the public query node $\tau$ the null picks $p\in[0,1/2]$; nodes with different $\tau$ are publicly distinguishable. A mixture with masses $a_\tau$ (both branches), $b_\tau$ ($Y=1$ only), $c_\tau$ ($Y=0$ only) and pre-release mass $m$ therefore has adaptive risk $\sum_\tau[a_\tau+\max(c_\tau,(b_\tau+c_\tau)/2)]+m$, because the null maximises $(1-p)Z_0+pZ_1$ at an endpoint, and i.i.d.\ risk $\sum_\tau[a_\tau+(b_\tau+c_\tau)/2]+m$. Each problem maximises a linear gain under a sum of positively homogeneous convex gauges, so its optimum is $\alpha$ times the best gain per unit risk over the extreme rays. Under adaptive risk the rays are $b$ alone (ratio $2g_{z1}$), $c$ alone ($g_{z0}$), $c=b$ (ratio at most $g_{\mathrm{both}}$, dominated), $a$ ($g_{\mathrm{both}}$) and pre-release ($g_{\mathrm{pre}}$). Under i.i.d.\ risk, $c$ alone has ratio $2g_{z0}$. Since $\alpha\le1/2$, the risk-$1/2$ rays are feasible at mass $2\alpha\le1$.
\end{proof}

\paragraph{The symmetric chain.} Take $\eta_{01}=\eta_{10}=\eta\in(1/2,1)$, $\rho=1-2\eta=-r$, prior $1/2$, emissions $(e,1-e)$ with $a=1-2e$ and $g=a^2$.

\begin{lemma}
\label{lem:S}
(i) $\PP_A(Y_\tau=1)=1/2$ for every $\tau$. (ii) Before any label is released, $q_t=1/2$. (iii) After $Y_\tau=y$ is released, the gain of the challenger over the incumbent at time $\tau+j$, $j\ge d$, is $w_y(j)=2q_{\tau+j}(y)-1=(2y-1)\,g\,\rho^j$.
\end{lemma}

\begin{proof}
$\PP(S_t=1)=1/2$ for all $t$, so $\PP(Y_t=1)=1/2$, which gives (i) and (ii). Bayes gives $\PP(S_\tau=1\mid Y_\tau=y)=1/2+(2y-1)a/2$; the symmetric kernel propagates deviations from $1/2$ by the factor $\rho$ per step; and $2q-1=a\cdot2(\PP(S=1)-1/2)$. A label released at $\tau+d$ conditions the action at $\tau+d$, so $j\ge d$.
\end{proof}

So branch $y$'s positive steps are the offsets $j$ with $(-1)^j(2y-1)>0$: odd $j$ after $Y=0$ and even $j$ after $Y=1$. Let $G_y(\tau)$ be the largest value of $\sum_jx_jw_y(j)-k\,\mathrm{sw}(x)$ over nonzero 0/1 schedules $x$ on $j\in[d,T-\tau]$, and $\bar G_y=\max_\tau G_y(\tau)$. By Lemma~\ref{lem:S}(i), $g_{zy}(\tau)=G_y(\tau)/2-c$ and $g_{\mathrm{both}}(\tau)=(G_0(\tau)+G_1(\tau))/2-c$.

\begin{lemma}
\label{lem:pre}
In the symmetric chain $g_{\mathrm{pre}}\le\max(0,\max_\tau\max\{g_{z0}(\tau),g_{z1}(\tau),g_{\mathrm{both}}(\tau)\})$.
\end{lemma}

\begin{proof}
Take a plan that promotes unconditionally. If no label is released by $T$, all gains are 0 by Lemma~\ref{lem:S}(ii), so the plan's reduction is at most 0. If a label is queried at $\tau$ and released at $\tau+d\le T$, zero out every action before $\tau+d$ in each branch: those gains are 0, and the number of switches cannot rise. Each branch is then empty or a post-release schedule, so the reduction is at most that of the design with promotion set $\{y:\text{branch nonempty}\}$, i.e.\ $g_{\mathrm{both}}$, $g_{z0}$, $g_{z1}$ or $-c$.
\end{proof}

\paragraph{Proof of Theorem~\ref{thm:premium}(1).} Let $c=k=0$, $d$ odd and $n=T-d\ge1$. With $k=0$, covering every positive step is optimal and $G_y(\tau)$ is nondecreasing in the window length, so $\bar G_0$ and $\bar G_1$ are both attained at $\tau=1$ and $\max_\tau(G_0(\tau)^++G_1(\tau)^+)=\bar G_0+\bar G_1$. Since $d$ is odd, branch 0 covers the $n_0=\lceil n/2\rceil$ odd offsets in $[d,T-1]$ and branch 1 the $n_1=\lfloor n/2\rfloor$ even ones:
\[
\bar G_0=gr^d\frac{1-r^{2n_0}}{1-r^2},\qquad \bar G_1^+=gr^{d+1}\frac{1-r^{2n_1}}{1-r^2}\le\bar G_0,
\]
where for $n=1$ branch 1 has no positive step and its best design is worth less than 0, so only $\bar G_1^+=0$ enters below. Write $\bar G_1$ for $\bar G_1^+$ from here on.
By Lemmas~\ref{lem:N1} and~\ref{lem:pre}, $r_c=(\bar G_0+\bar G_1)/2$ and $r_{\mathrm{IID}}=\bar G_0$, so $\Gamma=\alpha(\bar G_0-\bar G_1)/2$ and $\mathrm{VOI}_{\mathrm{IID}}=\alpha\bar G_0$ for every $\alpha\le1/2$. Hence
\[
\rhoG(T)=\frac{\bar G_0-\bar G_1}{2\bar G_0}=
\begin{cases}\dfrac{1-r}{2}, & n\text{ even},\\[2mm]
\dfrac{(1-r)(1+r^n)}{2(1-r^{n+1})}\ \ge\ \dfrac{1-r}{2}, & n\text{ odd},\end{cases}
\]
which does not involve $\alpha$, and $\inf_{T\ge d+1}\rhoG(T)=\lim_T\rhoG(T)=(1-r)/2$. \hfill$\square$

\subsection{The first-label relaxation, and Theorem~\ref{thm:premium}(2)}

A pure plan has a deterministic first query time $\tau$. Class $\mathsf P$: plans that promote before $\tau+d$ or never query; their events are constant. Class $\mathsf Q$: plans with no promotion before $\tau+d$. For $\pi\in\mathsf Q$ let $g_y(\pi)$ be the expected cost saving given $Y_\tau=y$ (all costs included) and $r_y(\pi)=\PP_{\mathrm{IID}}(\text{promote}\mid Y_\tau=y)$. In the symmetric chain $\PP_A(Y_\tau=y)=1/2$, so $\mathrm{gain}(\pi)=(g_0+g_1)/2$ and $R_{\mathrm{IID}}(\pi)=(r_0+r_1)/2$. Let $H_\varnothing=\max_{\pi\in\mathsf P}\mathrm{gain}(\pi)$ and $s_y=\max\{g_y(\pi)/r_y(\pi):\pi\in\mathsf Q,\ r_y(\pi)>0\}$. A branch with $r_y=0$ never promotes and only pays costs, so $g_y(\pi)\le s_y^+r_y(\pi)$ for every $\pi\in\mathsf Q$.

\begin{lemma}[First-label relaxation]
\label{lem:M3}
In the symmetric chain, for every $T$, $B$, $d$, $c$, $k$ and $\alpha\in(0,1]$, with $S'=\max(H_\varnothing^+,s_1^+,(s_0^++s_1^+)/2)$ and $S_I=\max(H_\varnothing^+,s_0^+,s_1^+)$:
(i) $\mathrm{VOI}_{\mathrm{causal}}\le\alpha S'$; (ii) $\mathrm{VOI}_{\mathrm{IID}}\le\alpha S_I$; (iii) for any pure plan $\hat\pi$ with $R_{\mathrm{IID}}(\hat\pi)\ge\alpha$ and ratio $\zeta=\mathrm{gain}(\hat\pi)/R_{\mathrm{IID}}(\hat\pi)\ge0$, $\Gamma\ge\alpha(\zeta-S')$ and $\rhoG\ge(\zeta-S')/S_I$.
\end{lemma}

\begin{proof}
Let $\nu_1$ be the null with $p=0$ at the first query and $p=1/2$ afterwards; it is adaptive because the first query is public. For $\pi\in\mathsf Q$, $R_{\nu_1}(\pi)=r_0(\pi)$. (i) Let $\mu$ be causal-feasible with mass $m$ on the promoting plans of $\mathsf P$ (plans in $\mathsf P$ that never promote have risk 0 and gain at most 0, so dropping them can only increase the gain), and put $A=\sum_{\pi\in\mathsf Q}\mu(\pi)r_0(\pi)$, $D=\sum_{\pi\in\mathsf Q}\mu(\pi)r_1(\pi)$. Then $\mathrm{gain}(\mu)\le mH_\varnothing^++(s_0^+A+s_1^+D)/2$. If $D\le A$ this is at most $mH_\varnothing^++A(s_0^++s_1^+)/2$; if $D>A$ it is at most $mH_\varnothing^++S'A+S'(D-A)/2$. In both cases $\mathrm{gain}(\mu)\le S'(m+\max(A,(A+D)/2))=S'\max(R_{\nu_1}(\mu),R_{\mathrm{IID}}(\mu))\le S'\alpha$. (ii) The same computation with the i.i.d.\ row only. (iii) Mixing $\hat\pi$ with never-promote at probability $\alpha/R_{\mathrm{IID}}(\hat\pi)$ is i.i.d.-feasible with value $\alpha\zeta$; subtract (i) and divide by (ii).
\end{proof}

Now let $B=2$, $k=c=0$, $d$ odd, and put $O_d=r^d/(1-r^2)$, $\bar G_0^\infty=gO_d$, $\bar G_1^\infty=r\bar G_0^\infty$ and $v_1=gr^d/(2(1-r))$. With $k=0$ the post-promotion continuation plays $a_s=\one\{w_s>0\}$, so its value is $\sum_{s\ge\sigma}\E(w_s)^+$, nondecreasing in $T$. For a state deviation $x$ at the query time $u$, $\PP(Y_u=1)=(1+ax)/2$ and the posterior deviation is $(x\pm a)/(1\pm ax)$, so $\PP(y)x'_y=(x+\epsilon_ya)/2$ with $\epsilon_y=2y-1$. By positive homogeneity of $z\mapsto z^+$, every expected post-label value is a sum of positive parts of linear combinations of powers of $\rho$.

\begin{lemma}[$T=\infty$ branch values]
\label{lem:M41}
For a first label at $u$ and a second at $u+\delta$, $\delta\ge1$, and $y\in\{0,1\}$: promoting at $u+d$ and using the second label freely is worth $H_y(\delta)=\bar G_y^\infty+v_1-gr^{d+\delta}/(2(1-r))$; promoting only on the second label's value $v$ (branch risk $1/2$) has ratio $g(1+\epsilon_y\rho^\delta)E_d$ for $v=1$ and $g(1-\epsilon_y\rho^\delta)O_d$ for $v=0$, with $E_d=r^{d+1}/(1-r^2)$; and unconditional promotion with both labels free is worth $H_\varnothing\le2v_1$.
\end{lemma}

\begin{proof}
Given $Y_u=y$ the deviation at $u$ is $\epsilon_ya$. Before the second release the value is $\sum_{j=d}^{\delta+d-1}g(\epsilon_y\rho^j)^+$. After it, with $x=\epsilon_ya\rho^\delta$, the value at offset $j\ge d$ is $\tfrac g2[((\epsilon_y\rho^\delta+1)\rho^j)^++((\epsilon_y\rho^\delta-1)\rho^j)^+]=\tfrac g2[(1+\epsilon_y\rho^\delta)r^j\one\{j\text{ even}\}+(1-\epsilon_y\rho^\delta)r^j\one\{j\text{ odd}\}]$, using $|\rho^\delta|<1$. Summing gives $v_1-\epsilon_y\tfrac{gr^d}{2}\tfrac{\rho^\delta}{1+r}$. Adding the first segment and splitting $\delta$ by parity gives, in all four cases, $\bar G_y^\infty+v_1-gr^{d+\delta}/(2(1-r))$. The pattern ratios follow from $\PP(v\mid y)x''=(\epsilon_ya\rho^\delta+\epsilon_va)/2$. For $H_\varnothing$, average the immediate values over $y$: $(H_0+H_1)/2<(\bar G_0^\infty+\bar G_1^\infty)/2+v_1=2v_1$.
\end{proof}

\begin{lemma}[Branch slopes at $B=2$]
\label{lem:M42}
For every $T$, $s_0\le s_0^\infty:=\bar G_0^\infty+v_1=gr^d(3+r)/(2(1-r^2))$, $s_1\le s_1^\infty:=gr^d/(1-r)=2v_1$ and $H_\varnothing\le2v_1$.
\end{lemma}

\begin{proof}
In branch $y$ one label remains, so the branch event is empty, everything, or $\{Y_2=v\}$. ``Everything'' is worth at most $\sup_\delta H_y(\delta)=\bar G_y^\infty+v_1$; $\{Y_2=v\}$ has branch risk $1/2$ and ratio at most $\sup_{\delta,v}$ of the pattern ratios. For $y=0$ the best pattern is $\delta=2$, $v=0$, worth $g(1+r^2)O_d$, and $(3+r)/2-(1+r^2)=(1-r)(1+2r)/2>0$, so $s_0\le\bar G_0^\infty+v_1$. For $y=1$ the best pattern is $\delta=1$, $v=0$, worth $g(1+r)O_d=gr^d/(1-r)$, which exceeds $\bar G_1^\infty+v_1=gr^d(1+3r)/(2(1-r^2))$. A first query later than step 1 only shortens the horizon, by stationarity.
\end{proof}

\paragraph{Proof of Theorem~\ref{thm:premium}(2).} Let $\zeta(T)=\max_\delta H_0(\delta;T)$ be the finite-horizon value of the plan ``query at 1; if $Y_1=0$ promote at $1+d$, query again at $1+\delta$ and follow the sign of the gain; if $Y_1=1$ do nothing''. It has i.i.d.\ risk $1/2\ge\alpha$ and ratio $\zeta(T)$. Lemma~\ref{lem:M42} gives $S'\le\max(2v_1,s_1^\infty,(s_0^\infty+s_1^\infty)/2)=(s_0^\infty+s_1^\infty)/2$ and $S_I\le s_0^\infty$, so Lemma~\ref{lem:M3}(iii) yields
\[
\rhoG(T)\ \ge\ \frac{[\zeta(T)-(s_0^\infty+s_1^\infty)/2]^+}{s_0^\infty}.
\tag{D.1}
\]
Truncating at $T$ loses at most $\sum_{j\ge T-\delta}gr^j$ after the second release, so $\zeta(T)\ge s_0^\infty-g(r^{d+\delta}/2+r^{T-\delta})/(1-r)$ for every $1\le\delta\le T-d$, and $\zeta(T)\to s_0^\infty$. Since $s_0^\infty-s_1^\infty=gr^d/(2(1+r))$,
\[
\liminf_T\rhoG(T)\ \ge\ \frac{s_0^\infty-s_1^\infty}{2s_0^\infty}=\frac{1-r}{2(3+r)}.\qquad\square
\]

\subsection{Fixed budgets: Theorem~\ref{thm:premium}(4) and (3)}

Here $d=1$, $c=k=0$, $\alpha=1/10$. Let $X_t\in\{-1,1\}$ be the spin of the symmetric chain, $\E[X_{t+1}\mid X_t]=\rho X_t$, so the challenger's gain is $a\,\E[X_t\mid\text{released labels}]$; a query sign $S_t\in\{-1,1\}$ has $\PP(S_t=s\mid X_t)=(1+asX_t)/2$, and a zero label is $s=-1$. Fix $B\ge1$, put $H_B=6B^3+3B+2$, and assume $2/5<g<1/2$ and $0<r<r_0(B)=1/(16(B+2)H_B)$.

\paragraph{Step 1: cylinder densities.} For ordered times $v_1<\dots<v_{2q}$, $\E\prod_\ell X_{v_\ell}=\rho^{\sum_{\ell=1}^q(v_{2\ell}-v_{2\ell-1})}$ and odd moments vanish (write each $X_v$ as the initial spin times the transition signs up to $v$). Fix the learner's private seed and a realised decision-tree cylinder with $m\le B$ released signs $s_1,\dots,s_m$ at query times $u_1<\dots<u_m$; these times are fixed on the cylinder because each query decision is determined by earlier signs. Expanding $\E\prod_i(1+as_iX_{u_i})$ gives alternative probability $2^{-m}D_m$ with
\[
D_m=\sum_{J\subseteq[m],\,|J|\text{ even}}a^{|J|}s_J\,\rho^{\sum_\ell(u_{j_{2\ell}}-u_{j_{2\ell-1}})},
\tag{D.2}
\]
and, for $n\ge u_m+1$ with no further labels, posterior gain
\[
w_n=\frac1{D_m}\sum_{J\subseteq[m],\,|J|\text{ odd}}a^{|J|+1}s_J\,\rho^{\sum_{\ell\le(|J|-1)/2}(u_{j_{2\ell}}-u_{j_{2\ell-1}})+n-u_{j_{|J|}}}.
\tag{D.3}
\]

\paragraph{Step 2: uniform bounds.} Put $x=B^2r<1/288$. Every paired gap is at least one and $g<1$, so $|D_m-1|\le\delta_B(r):=\sum_{q\ge1}\binom{B}{2q}r^q\le x/(1-x)\le2B^2r$. In (D.3) isolate the latest singleton, whose term is $gs_m\rho^{n-u_m}$; its positive-part tail is $G_0=gr/(1-r^2)$ after a zero label and $G_1=gr^2/(1-r^2)$ after a one. Older singletons have $n-u_i\ge2$, and an odd subset of size $2q+1$ carries at least $q$ powers of $r$ from its paired gaps. Summing absolute values over $n>u_m$ and dividing by $gr$ bounds the rest by $P_B(r)\le2Br+4B^3r$.

\paragraph{Step 3: Bellman envelope.} For a spin mean $x$, the posterior after sign $s$ is $\psi_s(x)=(x+as)/(1+asx)$ with probability $(1+asx)/2$; then $\E\psi_S(x)=x$, $\E|\psi_S(x)|=\max(|x|,a)$, and for $j\ge1$
\[
\E[(a\rho^j\psi_S(x))^+]-(a\rho^jx)^+=\tfrac{gr^j}{2}\big(1-|x|/a\big)^+.
\tag{D.4}
\]
Let $F_t(x)=\sum_{j=0}^{T-t}(a\rho^jx)^+$ and $W_{t,b}(x)$ the largest expected future gain with $b$ queries left. Backward induction over $(t,b)$ gives $W_{t,b}(x)\le F_t(x)+bv_1$ with $v_1=gr/(2(1-r))$: the current action contributes at most $(ax)^+$; without a query the continuation is at most $F_{t+1}(\rho x)+bv_1$; with a query (released at $t+1$) it is at most $\E F_{t+1}(\rho\psi_S(x))+(b-1)v_1$, and by (D.4) the expected increase over $F_{t+1}(\rho x)$ is at most $\sum_{j\ge1}gr^j/2\le v_1$. Since actions do not affect observations and costs are zero, this covers rollback, re-promotion and a query on the promotion step.

\paragraph{Step 4: leaf bounds.} On a first-promotion leaf $h$ with $m\ge1$ labels, latest label $y$ and $n=B-m$ queries left, Steps 2 and 3 give $D_mV(h)/(gr)\le G_y/(gr)+P_B(r)+(1+\delta_B)n/(2(1-r))$. For $r<1/2$, $G_0/(gr)-1\le2r$, $G_1/(gr)\le2r$ and $\tfrac n2(\tfrac{1+\delta_B}{1-r}-1)\le2B^3r+Br$, so
\[
D_mV(h)/(gr)\ \le\ \one\{y=0\}+n/2+H_Br.
\tag{D.5}
\]
For $m=0$ the spin mean is zero and $V(h)/(gr)\le B/(2(1-r))\le B/2+Br$.

\paragraph{Step 5: null certificate.} Let $\nu_j$ force the $j$th queried label to zero and use fair labels elsewhere; it is adaptive because the query ordinal is public. The mixture $Q=(\mathrm{IID}+\nu_1+\dots+\nu_B)/(B+1)$ gives a cylinder with $z$ zeros the mass $Q(h)=2^{-m}(B+1-m+2z)/(B+1)$, so $\PP_A(h)/Q(h)=(B+1)D_m/(B+1-m+2z)$. First-promotion leaves are disjoint and gain before promotion is zero, so for each seed $\mathrm{gain}\le R_Q\max_h\PP_A(h)V(h)/Q(h)$. If the latest label is zero then $z\ge1$ and the leading term is at most $\frac{B+1}{n+3}(1+\frac n2)\le\frac{(B+1)^2}{2(B+2)}=:C_B$; if it is one, at most $\frac{B+1}{n+1}\frac n2\le C_B$; for $m=0$ it is $B/2\le C_B$. The prefactor on the remainder is at most $B+1$. Averaging over the seed (no conditional feasibility is needed, since $R_Q$ is an average of admissible nulls' risks),
\[
\mathrm{VOI}_{\mathrm{causal}}\le\alpha gr[C_B+(B+1)H_Br],\qquad \mathrm{VOI}_{\mathrm{IID}}\le\alpha gr[I_B+H_Br],\quad I_B=\tfrac{B+1}{2}.
\tag{D.6}
\]

\paragraph{Step 6: i.i.d.\ witness.} For $B\ge2$: query at 1; promote at 2 exactly when $Y_1=0$; on that branch query at $5,8,\dots,3B-1$ and play the sign of the posterior gain. Its i.i.d.\ risk is $1/2$. The time-2 gain is $gr$; the time-5 query has prior mean of magnitude $ar^4$ and, by (D.4), adds $gr(1-r^4)/2$; each later query has prior magnitude below $a$ with $\E[|x|/a]=r^3$ and adds $gr(1-r^3)/2$. Release times are distinct and $T\ge3B$ includes them all, so $\zeta_T/(gr)\ge\frac{B+1}{2}-\frac{r^4+(B-2)r^3}{2}$; for $B=1$, $\zeta_T/(gr)\ge1=I_1$. Mixing with never-promote to risk $\alpha$ gives $\mathrm{VOI}_{\mathrm{IID}}\ge\alpha\zeta_T$.

\paragraph{Proof of Theorem~\ref{thm:premium}(4).} The witness remainder and the certificate remainder are each below $(B+1)/(16(B+2))$ under $r<r_0(B)$, and $I_B-C_B=(B+1)/(2(B+2))$. Hence $\Gamma\ge\alpha gr\,\frac{3(B+1)}{8(B+2)}$. Also $I_B+H_Br<B+1$, and dividing by (D.6) gives $\rhoG(T)\ge3/(8(B+2))$ for every $T\ge3B$. Since $H_B\le11B^3$ and $B+2\le3B$, $r_0(B)\ge1/(528B^4)$: for $0<r<1/1056$ every $B\le\lfloor(1056r)^{-1/4}\rfloor$ is covered. \hfill$\square$

\paragraph{Proof of Theorem~\ref{thm:premium}(3).} On the rectangle $r\in(1/200,1/100)$, $g\in(2/5,1/2)$:
\begin{itemize}
\item $B=1$: part 1 gives $\rhoG\ge(1-r)/2>99/200$.
\item $B=2$: take $\delta=2$ in (D.1) with the truncation bound; with $d=1$ this gives
\[
\rhoG\ \ge\ \Big[\tfrac1{4(1+r)}-\tfrac{r^2/2+r^{T-3}}{1-r}\Big]\Big/\tfrac{3+r}{2(1-r^2)},
\]
whose numerator decreases and denominator increases in $r$, and which increases in $T$; at $r=1/100$, $T=20$ it exceeds $0.16$.
\item $B=3,\dots,6$: Steps 1--6 with the remainders kept term by term instead of bounded by $P_B$ and $\delta_B$. For $B=3$ the remainder after the latest singleton is $g[s_1\rho^{n-u_1}+s_2\rho^{n-u_2}]+g^2s_1s_2s_3\rho^{u_2-u_1+n-u_3}$, giving $E=[G_0g(2r+r^2)+g(r^2+r^3)/(1-r)+g^2r^2/(1-r)]/(1-g(2r+r^2))$ in place of $H_Br$; $B=4,5,6$ retain the triple, quintuple and six-spin terms in the same way. The leaf maxima of (D.5)--(D.6) are then evaluated by exact rational enumeration of all histories at the closed corner $(r,g)=(1/100,1/2)$, which is valid on the open rectangle because every normalised envelope term increases in $r$ and $g$. The resulting coefficients are, for $B=3,4,5,6$: causal $5360/3333$, $62875/29997$, $60400/23331$, $41125/13332$; i.i.d.\ $6700/3333$, $25150/9999$, $30200/9999$, $11750/3333$. The witnesses of Step 6 give $\zeta_T/(gr)>1.99,\,2.49,\,2.99,\,3.49$ for $T\ge16$ ($T\ge20$ for $B=6$). Hence $\rhoG\ge19/101$, $13/84$, $40/303$, $13/118$.
\end{itemize}
The minimum over $B\le6$ exceeds $0.10$ for every $T\ge20$. \hfill$\square$

The enumeration for part 3 is computer-assisted in exact rational arithmetic; the scripts are in the supplement. Parts 1, 2 and 3 were also checked numerically against certified exact values (Section~\ref{sec:e3}).
